\documentclass{article}
\usepackage[T1]{fontenc}
\usepackage{lmodern}
\usepackage{graphicx}
\usepackage{amsmath,amssymb}
\usepackage[a4paper,margin=2cm]{geometry}
\usepackage[absolute,overlay]{textpos}
\setlength{\TPHorizModule}{1mm}
\setlength{\TPVertModule}{1mm}
\usepackage[indonesian]{babel}
\usepackage{hyperref}
\setlength{\emergencystretch}{2em}
% unit: Halaman 13

\begin{document}

% === Halaman 13 (llm) ===

\section*{2. Blind stegananalysis}

\begin{itemize}
    \item Teknik steganalisis yang bekerja pada sembarang algoritma steganografi dan sembarang format media.
    \item Teknik ini mempelajari perbedaan antara statistik \textit{cover-object} dan \textit{stego-object} dan membedakannya. Proses pembelajaran (\textit{learning}) dilakukan dengan melatih (\textit{training}) mesin pada sekumpulan database media. Model \textit{machine learning} yang digunakan misalnya jaringan syaraf tiruan.
    \item Hasil steganalisis kurang akurat dibandingkan dengan teknik \textit{targeted steganalysis}, tetapi kelebihannya adalah dapat diperluas untuk algoritma yang lain.
\end{itemize}

\end{document}
